\scp{Summary and implications}{15-07}
In the preceding chapters we have explored a number of typological
and syntactic features found unevenly in some languages \ili{east}
of the Blust line,
but not west of it.
Many of these features are found in the typologically SOV Papuan
languages still spoken in the region,
and are attributed to contact rather than to inheritance (see \srf{15-Features}).
We have looked carefully at the historical phonologies
and morphosyntax of Austronesian languages in the region,
and have identified nine separate subgroups,
one of which links to \tsc{Celebic} languages (see chapters \ref{ch:BL}--\ref{ch:Ban}).
We are unable to link the remaining eight Wallacean subgroups
together in the historical phonology at any identifiable node below PMP.
We explored various subgrouping arguments that have been put
forward to link some or all of these subgroups together,
and have been unable to find any that withstand scrutiny (see chapters \ref{ch:OthSub}--\ref{ch:Ban}).

Some arguments in support of CEMP have been based on the lexicon.
We observed that the lexicon is the weakest
and most problematic basis for subgrouping,
because words are easily borrowed.
In \srf{15-04} we looked at a number of words
and features, and their distribution, found in both Austronesian
and Papuan languages (spread through contact); words
and features that straddle both sides of the Blust line but cannot
be assigned to PMP (less than MP,
but greater than CEMP); words
and features that are found in two or more Wallacean subgroups,
but not in all of them (less than CMP); words
and features that are found unevenly in some nodes of several
subgroups, but not in all of them (spread by contact
and diffusion, cannot be assigned to the parent language); words
and features that are found unevenly \ili{east} of the Blust line,
but not west of it (almost CEMP).
The uneven and limited distribution of all of these categories
of words, often interspersed with other words having the same meaning,
strongly suggests they spread through contact
and diffusion and not through inheritance.
This should not be surprising in a region of high mobility
and at l\ili{east} two millennia of significant trade.
We have named a number of regional interlanguages (languages
of wider communication) that have been documented over the past
500 years, and suspect there are others that have come
and gone of which we are not aware,
and noted that these interlanguages may be partly responsible
for the uneven distribution of so many lexical items (\srf{15-05}).

Words for banana shown in \frf{15-11} are in some ways a representative
visual reflection of the complexities we find in Linguistic Wallacea.\footnote{
		\frf{02-20}
		focused only on reflexes of {\#}muku.
		This figure adds a number of other etyma to show the variety
		and complexity of associated terms in the region.}
Note in \frf{15-11}
that: several etyma show Papuan-Austronesian crossover requiring
spread by contact; some terms have Austronesian origins,
some have Papuan origins,
and we don't know about the others; all etyma (except North Halmahera
*b[ao]le) show uneven distribution with pockets of concentration
and pockets of absence; several regions (such as Sumba,
eastern Timor, \ili{Aru}, \tsc{\ili{Ambon}-Seram}, and the Bird’s Neck) show
two or more etyma interspersed side-by-side within the same subgroups
and nodes of subgroups indicating spread by diffusion; some regions
(such as \tsc{Sula-Buru} and \tsc{Taliabu}) have neither retentions of PMP *punti,
nor of Papuan-sourced {\#}muku; reflexes of both *punti
and **kalauR straddle the Blust line.
Again, \frf{15-11} is presented here as a single visual representation
reflecting the variety and complexities found in many other words
and features in the historical phonologies,
morphosyntax and lexicon of the region,
as we have shown in this study.

\begin{figure}[h]\centering
	\includegraphics[width=\columnwidth]{pisang.pdf}
	\fcp{Distribution of several etyma meaning ‘banana’}{15-11}
\end{figure}

If we put a number of features that have been associated with
Wallacean languages together in a table,
it becomes clear that almost every language has a different profile.
One feature does not appear to align with others.
For example, \trf{15-09} shows a number of structural features
for select languages, including outside the target region,
along with their value for subgrouping,
and \trf{15-10} shows the presence of reflexes of selected words
in the lexicon in some Wallacean languages.
(Note: (+) indicates a feature does not occur in the same way
as it does for other languages; ? means we have insufficient
information; N/A = not applicable, i.e.
outside of range of marsupials.)

Some comments on the human/non-human distinction in \trf{15-09}
in \ili{Luang}, \ili{Selaru}, \ili{Dobel} and SW \ili{Tarangan} are needed.
In \ili{Luang}, non-human markers are a set of spatial,
temporal and referential deictic markers that can stand as the
head of a phrase \citep[39f, 83f]{TaberTaber2015}.
In \ili{Selaru}, non-human plural Actors cannot use a pronoun \citep[35]{Coward2005},
and is “{\0} for non-human/inanimate possessors” in the genitive construction (p.60).
These are quite different issues
and mechanisms from the human/non-human distinctions in \tsc{\ili{Ambon}-Seram} and SHWNG languages.
The \ili{Dobel} distinction is described as animate/inanimate.
The distinction is mostly limited to \tsc{3sg} Undergoer enclitics \citep[140]{Hughes2000}.
Free pronouns and Actor proclitics are the same for animate and inanimate.
Numerals are marked differently for animate and inanimate referents \citep[155--157]{Hughes2000}.
It is possible that SW. \ili{Tarangan} \it{nɛy} ‘\tsc{3sg}’ refers only
to humans (and God, and maybe other spirits),
but that is not certain.

Similarly, some comments on the complex numerals in \trf{15-09} are needed.
\ili{Malay}, \ili{Buru}, \ili{Sanana}, and \ili{Soboyo} have subtractive complex numerals,
while those in S. \ili{Mambae}, \ili{Dobel}, SW.
\ili{Tarangan}, and \ili{Kowiai} are additive and/or multiplicative (e.g. 6+1).
\ili{Laboya} has one multiplicative numeral (\it{podo pata} ‘eight’
= 2×4) and one subtractive numeral (\it{kaɓandiha} ‘nine’ = 10-1).

\begin{table}[p]\centering\tcp{Select structural features in Wallacean languages}{15-09}
%\begin{adjustbox}{width=\textwidth}
	\begin{threeparttable}\st{2pt}
	\begin{tabular}{ll|cccccc}\lsptoprule
Feature	&				&	Pre-posed &	Alienable/	&	Human/		&	Clause-							&	Complex	&	no ŋ	\\
				&				&	Poss.			&	Inalienable	&	Non-human	&	final \tsc{neg}			&	numbers	&		\\
\mc{2}{l|}{\it{Subgroup valueʔ}}		&	\it{contact}	&	\it{contact}	&	\it{contact}	&	\it{contact}	&	\it{contact}	&	\it{areal}	\\	\midrule
\ili{Malay} 	&			&	--	&	--	&	--	&	--	&	+ {\cdr}	&	--	\\
\ili{Cia-cia}	&	\tsc{Cel}	&	--	&	--	&	?	&	--	&	--	&	--	\\
\ili{Bima}	&	BL	&	--	&	--	&	?	&	--	&	--	&	--	\\
\ili{Manggarai} 	&	BL	&	--	&	--	&	--	&	--	&	--	&	--	\\
\ili{Kambera}	&	BL	&	--	&	--	&	--	&	--	&	--	&	--	\\
\ili{Laboya}	&	BL	&	--	&	--	&	--	&	\hp{\su{†}}(+)\su{†} {\cdr}	&	+ {\cdr}	&	--	\\
\ili{Havu} 	&	BL	&	--	&	--	&	--	&	--	&	--	&	--	\\
\ili{Helong}	&	TB	&	+ {\cdr}	&	+ {\cdr}	&	--	&	+ {\cdr}	&	--	&	--	\\
\ili{Tii} 	&	TB	&	+ {\cdr}	&	--	&	--	&	--	&	--	&	--	\\
K. \ili{Amarasi}	&	TB	&	+ {\cdr}	&	--	&	--	&	\hp{\su{†}}(+)\su{†} {\cdr}	&	--	&	+ {\cdr}	\\
\ili{Welaun}	&	CT	&	+ {\cdr}	&	+ {\cdr}	&	?	&	--	&	--	&	+ {\cdr}	\\
\ili{Kemak}	&	CT	&	+ {\cdr}	&	+ {\cdr}	&	--	&	--	&	--	&	+ {\cdr}	\\
S. \ili{Mambae}	&	CT	&	+ {\cdr}	&	--	&	--	&	--	&	+ {\cdr}	&	+ {\cdr}	\\
\ili{Tetun}	&	TB	&	+ {\cdr}	&	--	&	--	&	--	&	--	&	\hp{\su{‡}}+\su{‡} {\cdr}	\\
\ili{Galolen}	&	TB	&	+ {\cdr}	&	--	&	?	&	+ {\cdr}	&	--	&	+ {\cdr}	\\
\ili{Luang}	&	TB	&	+ {\cdr}	&	--	&	(+) {\cdr}	&	--	&	--	&	+ {\cdr}	\\
\ili{Selaru}	&	TB	&	+ {\cdr}	&	+ {\cdr}	&	(+) {\cdr}	&	--	&	--	&	+ {\cdr}	\\
\ili{Dobel}	&	\tsc{Aru}	&	+ {\cdr}	&	+ {\cdr}	&	+ {\cdr}	&	--	&	+ {\cdr}	&	--	\\
SW. Tarang.	&	\tsc{Aru}	&	+ {\cdr}	&	+ {\cdr}	&	(+) {\cdr}	&	--	&	+ {\cdr}	&	--	\\
\ili{Yamdena}	&	STB	&	+ {\cdr}	&	+ {\cdr}	&	--	&	--	&	--	&	--	\\
\ili{Kei} Kecil	&	STB	&	+ {\cdr}	&	+ {\cdr}	&	?	&	+ {\cdr}	&	--	&	--	\\
\ili{Geser}	&	STB	&	+ {\cdr}	&	--	&	?	&	+ {\cdr}	&	--	&	--	\\
\ili{Kowiai}	&	STB	&	+ {\cdr}	&	+ {\cdr}	&	?	&	+ {\cdr}	&	+ {\cdr}	&	--	\\
S. \ili{Nuaulu}	&	AS	&	+ {\cdr}	&	+ {\cdr}	&	+ {\cdr}	&	+ {\cdr}	&	--	&	\hp{\su{‡}}+\su{‡} {\cdr}	\\
\ili{Larike}	&	AS	&	+ {\cdr}	&	+ {\cdr}	&	+ {\cdr}	&	--	&	--	&	\hp{\su{‡}}+\su{‡} {\cdr}	\\
\ili{Asilulu}	&	AS	&	+ {\cdr}	&	+ {\cdr}	&	+ {\cdr}	&	--	&	--	&	\hp{\su{‡}}+\su{‡} {\cdr}	\\
\ili{Buru}	&	SB	&	+ {\cdr}	&	--	&	--	&	+ {\cdr}	&	+ {\cdr}	&	--	\\
\ili{Sanana}	&	SB	&	+ {\cdr}	&	--	&	\hp{\su{Δ}	}--\su{Δ}	&	+ {\cdr}	&	+ {\cdr}	&	--	\\
\ili{Soboyo}	&	\tsc{Cel}	&	+ {\cdr}	&	--	&	--	&	+ {\cdr}	&	+ {\cdr}	&	--	\\
\ili{Banggai}	&	\tsc{Cel}	&	+ {\cdr}	&	--	&	--	&	--	&	--	&	--	\\
\lspbottomrule
	\end{tabular}
		\begin{tablenotes}\tnsize
			\item[†]	{\ili{Laboya}
								and \ili{Kotos} \ili{Amarasi} have double negation.}
			\item[‡]	{Engma /ŋ/ is not in native vocabulary,
								but is found in a large number of loanwords.}
			\item[Δ]	{The putative non-human form \it{i-} is identical to the human
								form \it{i-} \citep[233, 417, 568]{Bloyd2020}.}
		\end{tablenotes}
	\end{threeparttable}
%\end{adjustbox}
\end{table}

\begin{table}[p]\centering\tcp{Reflexes or glosses of select words in Wallacean languages}{15-10}
\begin{adjustbox}{width=\textwidth}
	\begin{threeparttable}\st{2pt}
	\begin{tabular}{ll|ccccccc}\lsptoprule
Feature	&			&	**mantəR	&	{\#}muku	&	**malip	&	**maya	&	*duRi	&	*daləm	&	head leaf	\\
	&		&	cuscus	&	banana	&	laugh	&	tongue	&	bone	&	feelings	&	hair	\\
\mc{2}{l|}{\it{Subgroup value?}}			&	\it{no}	&	\it{contact}	&	\it{poss.}	&	\it{no}	&	\it{no}	&	\it{contact}	&	\it{contact}	\\	\midrule
\ili{Malay} 	&		&	N/A	&	--	&	--	&	--	&	--	&	--	&	--	\\
\ili{Cia-cia} 	&	\tsc{Cel}	&	+ {\cdr}	&	--	&	--	&	--	&	--	&	?	&	--	\\
\ili{Bima} 	&	BL	&	N/A	&	--	&	--	&	--	&	--	&	--	&	--	\\
\ili{Manggarai}	&	BL	&	N/A	&	+ {\cdr}	&	--	&	--	&	--	&	--	&	--	\\
\ili{Havu}	&	BL	&	N/A	&	+ {\cdr}	&	+ {\cdr}	&	--	&	+ {\cdr}	&	\hp{\su{†}}+/--\su{†} {\cdr}	&	+ {\cdr}	\\
\ili{Kambera}	&	BL	&	N/A	&	\hp{\su{‡}}(+)\su{†} {\cdr}	&	--	&	--	&	+ {\cdr}	&	--	&	--	\\
\ili{Laboya}	&	BL	&	N/A	&	?	&	+ {\cdr}	&	--	&	+ {\cdr}	&	--	&	--	\\
\ili{Helong}	&	TB	&	--	&	--	&	+ {\cdr}	&	+ {\cdr}	&	--	&	+ {\cdr}	&	--	\\
\ili{Tii}	&	TB	&	N/A	&	--	&	--	&	+ {\cdr}	&	+ {\cdr}	&	+ {\cdr}	&	--	\\
K. \ili{Amarasi}	&	TB	&	--	&	--	&	+ {\cdr}	&	+ {\cdr}	&	+ {\cdr}	&	--	&	--	\\
\ili{Welaun}	&	CT	&	+ {\cdr}	&	--	&	+ {\cdr}	&	--	&	+ {\cdr}	&	+ {\cdr}	&	+ {\cdr}	\\
\ili{Kemak} 	&	CT	&	+ {\cdr}	&	+ {\cdr}	&	--	&	--	&	+ {\cdr}	&	+ {\cdr}	&	--	\\
S. \ili{Mambae} 	&	CT	&	+ {\cdr}	&	+ {\cdr}	&	--	&	--	&	+ {\cdr}	&	+ {\cdr}	&	--	\\
\ili{Tetun}	&	TB	&	+ {\cdr}	&	--	&	--	&	--	&	+ {\cdr}	&	+ {\cdr}	&	--	\\
\ili{Galolen}	&	TB	&	+ {\cdr}	&	+ {\cdr}	&	+ {\cdr}	&	--	&	+ {\cdr}	&	+ {\cdr}	&	+ {\cdr}	\\
\ili{Luang}	&	TB	&	?	&	--	&	+ {\cdr}	&	--	&	+ {\cdr}	&	+ {\cdr}	&	--	\\
\ili{Selaru} 	&	TB	&	+ {\cdr}	&	--	&	+ {\cdr}	&	+ {\cdr}	&	+ {\cdr}	&	+ {\cdr}	&	--	\\
\ili{Dobel}	&	\tsc{Aru}	&	--	&	--	&	+ {\cdr}	&	--	&	--	&	--	&	--	\\
SW. Tarang.	&	\tsc{Aru}	&	\hp{\su{Δ}}(+)\su{Δ} {\cdr}	&	--	&	+ {\cdr}	&	--	&	--	&	\hp{\su{‖}}--\su{‖}	&	 --	\\
\ili{Yamdena}	&	STB	&	+ {\cdr}	&	--	&	+ {\cdr}	&	--	&	+ {\cdr}	&	+ {\cdr}	&	--	\\
\ili{Kei} Kecil	&	STB	&	+ {\cdr}	&	+ {\cdr}	&	+ {\cdr}	&	--	&	+ {\cdr}	&	+ {\cdr}	&	--	\\
\ili{Geser}	&	STB	&	--	&	--	&	+ {\cdr}	&	--	&	+ {\cdr}	&	?	&	--	\\
\ili{Kowiai}	&	STB	&	--	&	--	&	+ {\cdr}	&	--	&	+ {\cdr}	&	\hp{\su{‖}}--\su{‖}	 {\cdr}	&	--	\\
S. \ili{Nuaulu}	&	AS	&	+ {\cdr}	&	--	&	+ {\cdr}	&	+ {\cdr}	&	+ {\cdr}	&	+ {\cdr}	&	--	\\
\ili{Asilulu}	&	AS	&	+ {\cdr}	&	--	&	--	&	--	&	+ {\cdr}	&	+ {\cdr}	&	--	\\
\ili{Buru}	&	SB	&	--	&	--	&	+ {\cdr}	&	+ {\cdr}	&	+ {\cdr}	&	+ {\cdr}	&	--	\\
\ili{Sanana}	&	SB	&	--	&	--	&	--	&	--	&	+ {\cdr}	&	+ {\cdr}	&	--	\\
\ili{Soboyo}	&	\tsc{Cel}	&	--	&	--	&	--	&	--	&	--	&	+ {\cdr}	&	--	\\
\ili{Banggai}	&	\tsc{Cel}	&	--	&	--	&	--	&	--	&	--	&	?	&	--	\\
		\lspbottomrule
	\end{tabular}
		\begin{tablenotes}\tnsize
			\item[†]	{\ili{Havu} uses both *qatay > \it{ade} ‘liver’ and *daləm > \it{ɗara} ‘insides’
								for ‘seat of emotions’.} %in different phrases.}% \it{ade} `liver' is more frequent.}
			\item[‡]	{\ili{Kambera} has \it{muku} for a specific kind of banana.}
			\item[Δ]	{SW. \ili{Tarangan} \it{man} ‘bandicoot’
								%is anomalous in form
								%and meaning with other \tsc{Aru} languages (\srf{14-04},
								%\srf{14-05}) and 
								is not from **mantəR (\srf{14-04})}
			\item[‖]	{\ili{Kowiai} and SW. \ili{Tarangan} use words for `inside' which is not from *daləm as
								`seat of emotions' in idiomatic collocations.
								Thus, it is conceptually similar to *daləm ‘inside’,
								but not cognate in form.}
		\end{tablenotes}
	\end{threeparttable}
\end{adjustbox}
\end{table}

Looking at \trf{15-09}
and \trf{15-10}, it is fair to ask,
“Where is the patterned clustering of features across columns
that might have subgrouping value?” The features are sporadic,
do not align with each other,
and are full of discontinuities.
The same is true of their historical phonologies.
To a large degree,
that sums up the linguistic complexities in the Austronesian
languages and subgroups of Wallacea.
A few select features are plotted together in \frf{15-12} to
show the geographical range of many different thresholds.
The dotted lines do not indicate a hard boundary,
but rather an approximate
and gradual \it{starting point} where one begins to encounter
the feature indicated when travelling from west to \ili{east} or north to south.
Distribution maps like Figures \ref{fig:15-09}
and \ref{fig:15-11} show there are also many other associated etyma with
the same glosses scattered through the region than just the **maya
‘tongue’ and {\#}muku ‘banana’ shown here.
In other words, there are pockets of exceptions \ili{east}
and south of the dotted lines.

\begin{figure}[h]\centering
	\includegraphics[width=\columnwidth]{isoglosses.pdf}
	\fcp{Isoglosses where select features begin to be found}{15-12}
\end{figure}

\ili{As} can be seen in \frf{15-12}: several features are also found
in parts of Sulawesi; whether looking at the south,
the centre, or the north, these isogloss boundaries can be hundreds
of kilometres (and dozens of languages) apart from each other;
no single isogloss quite matches or defines the Blust line (to
some degree those features that approximate the Blust line between
SE Sulawesi and Timor-Maluku are different from those that mirror the Blust
line near \ili{Sumbawa}).

\it{Methodologically} we have noted that:

\begin{itemize}
	\item	Subgroups should have multiple corroborating evidence of different kinds.
				Exclusively shared innovations in the historical phonology provide
				a strong basis for subgrouping.
	\item	Subgroups should not depend on disputed data,
				disputed glosses, and non-overlapping distribution.
	\item	Since words are easily borrowed,
				subgroups should not depend solely on a few problematic words in the lexicon.
	\item	Following the comparative method,
				if witnesses are found in WMP
				and CMP, then a historical reconstruction can be assigned to
				the parent language PMP.
				But if WMP does not exist, and CMP does not exist,
				and no witnesses are found in western Indonesia or the Philippines,
				then the distribution and newer understanding of the subgrouping
				based on historical phonology requires that hundreds of PCMP
				and PCEMP reconstructions need to be reevaluated.
				Purported PCMP and PCEMP reconstructions need to be reassessed
				item-by-item to evaluate their accuracy,
				their validity, and the node or scope at which they are relevant.
	\item	PWMP reconstructions for which we have shown reflexes in
				our target region also need to be reassessed,
				but for different reasons.
	\item	We have shown that there are some etyma which have reflexes
				found in several subgroups across parts of central Maluku or
				Wallacea, that may reveal patterns in the lexicon above single Wallacean
				subgroups, but which have no legitimate proto-language to which to assign
				them, and which do not play a role in subgrouping arguments.
\end{itemize}